\chapter{Realization}
\minitoc
\newpage

\setcounter{secnumdepth}{0} % Set the section counter to 0 so next section is not counted in toc
% ----------------------- Introduction ----------------------- %
\section{Introduction}
In this chapter, we will finally discuss the long awaited realization.
This is the implementation phase where all of the work done in the previous chapters comes into picture.
This section will {\bf heavily focus on DevOps} since our primary objective has been migrating the \glsxtrshort{ilg} application to a new scalable architecture.
We will start with the new architecture hardware and software setup, talk a little bit about monitoring the application.
Later on, we will present some of the difficulties we encountered and then list all of the technologies we've used for our work.

\setcounter{secnumdepth}{2} % Resume counting the sections for the toc with a depth of 2 (Sections and sub-sections)
% ----------------------- Hardware setup ----------------------- %
\section{Hardware setup}
Since we agreed to proceed with the self-managed Kubernetes deployment as mentioned in \textbf{Chapter 2: State of the art}, we must create our own cluster nodes or virtual machines.
The nodes in question can be separated into the core nodes that are create once, and the scalable nodes that can be created indefinitely whenever we need to scale.

\subsection{Cloud provider}
To provision the hardware resources that we need for the deployment, we're going  to use the IONOS Cloud provider.
\cite[Ionos is a web hosting company founded in Germany in 1988 and is currently owned by United Internet.]{ionos-wikipedia}
\begin{figure}[H]
    \centering
    \makebox[\textwidth]{\includegraphics[width=10cm]{src/assets/images/ionos-logo.png}}
    \caption{Logo of IONOS}
    \label{fig:logo-of-ionos-cloud}
\end{figure}

IONOS Cloud provides -through the Ionos cloudpanel service- an easy way to create and manage virtual machines, firewall rules, private networks and various other infrastructure components.
\begin{figure}[H]
    \centering
    \makebox[\textwidth]{\includegraphics[width=15cm]{src/assets/images/ionos-cloudpanel-dashboard.JPG}}
    \caption{Ionos cloudpanel dashboard}
    \label{fig:image-of-ionos-cloudpanel-dashboard}
\end{figure}

The two images below show the available \glsxtrshort{vm} sizes in the Ionos cloudpanel dashboard and their corresponding specifications.
We will refer to the sizes in the tables of the next section entitled \textbf{Kubernetes Cluster}.
\begin{figure}[H]
    \centering
    \makebox[\textwidth]{\includegraphics[width=\linewidth]{src/assets/images/vm-standard.jpg}}
    \caption{Standard VM sizes in IONOS}
    \label{fig:standard-vm-sizes}
\end{figure}
\begin{figure}[H]
    \centering
    \makebox[\textwidth]{\includegraphics[width=\linewidth]{src/assets/images/vm-ram-optimized.JPG}}
    \caption{RAM Optimized VM sizes in IONOS}
    \label{fig:ram-optimized-vm-sizes}
\end{figure}

\subsection{Kubernetes Cluster}
Using the cloudpanel dashboard, it is time that we setup our cluster.
The tables below show the virtual machines that we need to get started.

\begin{table}[H]
    \renewcommand{\arraystretch}{1.5} % Padding
    \caption{Core cluster nodes}
    \centering
    \medskip
    \begin{tabularx}{1\textwidth} {
            | >{\hsize=1\hsize\linewidth=\hsize\raggedright\arraybackslash}X
            | >{\hsize=1.8\hsize\linewidth=\hsize\raggedright\arraybackslash}X
            | >{\hsize=0.2\hsize\linewidth=\hsize\centering\arraybackslash}X |}
        \hline
        \rowcolor{primary} \textbf{Name} & \textbf{Services}                            & \textbf{Size} \\
        \hline
        \textbf{kube-controller}         & The main Kubernetes cluster controller       & XL            \\
        \hline
        \textbf{data-node}               & PostgresDB, ilg-data, ilg-api, ilg-scheduler & L             \\
        \hline
        \textbf{front-node}              & ilg-front                                    & S             \\
        \hline
        \textbf{staging-node}            & PostgresDB, all other ilg services           & L             \\
        \hline
    \end{tabularx}
\end{table}

\begin{table}[H]
    \renewcommand{\arraystretch}{1.5} % Padding
    \caption{Scalable cluster nodes}
    \centering
    \medskip
    \begin{tabularx}{1\textwidth} {
            | >{\hsize=1.2\hsize\linewidth=\hsize\raggedright\arraybackslash}X
            | >{\hsize=1.5\hsize\linewidth=\hsize\raggedright\arraybackslash}X
            | >{\hsize=0.3\hsize\linewidth=\hsize\centering\arraybackslash}X |}
        \hline
        \rowcolor{primary} \textbf{Name} & \textbf{Services} & \textbf{Size} \\
        \hline
        \textbf{scraper-node-n}          & ilg-scraper       & RL            \\
        \hline
        \textbf{automation-node-n}       & ilg-automation    & RXL           \\
        \hline
    \end{tabularx}
\end{table}

Please note that
\begin{itemize}
    \item What each of the services represent is already mentioned in \textbf{Chapter 3: Analysis and specification of requirements}
    \item The size column which stands for the specification of the virtual machine is described in the previous section.
    \item The number of servers for \textbf{scraper}, and \textbf{automation} nodes will be scaled manually depending on the number of clients whenever the need arises.
\end{itemize}

\subsection{Proxy Servers}
Proxies are used to ensure all requests to LinkedIn come from a single IP address in order to avoid having to write a verification code every time.
\begin{table}[H]
    \renewcommand{\arraystretch}{1.5} % Padding
    \caption{List of proxy servers}
    \centering
    \medskip
    \begin{tabularx}{1\textwidth} {
            | >{\hsize=1\hsize\linewidth=\hsize\raggedright\arraybackslash}X
            | >{\hsize=1.8\hsize\linewidth=\hsize\raggedright\arraybackslash}X
            | >{\hsize=0.2\hsize\linewidth=\hsize\centering\arraybackslash}X |}
        \hline
        \rowcolor{primary} \textbf{Name} & \textbf{Services}                   & \textbf{Size} \\
        \hline
        \textbf{usp-n}                   & Upstream proxy server running Squid & XS            \\
        \hline
    \end{tabularx}
\end{table}

\newpage
% ----------------------- Software setup ----------------------- %
\section{Software setup}
After setting up our hardware resources, we end up with raw Linux virtual machines with no configuration whatsoever.
We have to create and manage our own Kubernetes instance on it to proceed any further.
We will be using {\bf MicroK8s}.
But before that, we also need to setup some basic configuration.

\subsection{Manual configuration}
For a single virtual machine, we have to do the following setup:

\begin{itemize}
    \item Create a user with sudo permissions
    \item Install the base packages we will be using
    \item Edit the hosts file so that the virtual machine recognizes the other virtual machines by their hostnames
    \item Install MicroK8s to have a single-node cluster running on the virtual machine
    \item Join the virtual machine with the master node to have a multi-node cluster.
\end{itemize}

\subsection{Automating the process}
Manually configuring the above can be extremely tedious, inefficient and error prone.
Therefore, we've created {\bf Ansible} scripts or playbooks that do just that. In the end, we simply edit a configuration file and directly run the script.
This has a huge advantage especially when
\begin{itemize}
    \item We want to create a development environment for easy prototyping.
    \item We want to use the exact same setup for future projects.
    \item We want to add nodes to the cluster later on when scaling.
\end{itemize}

The playbook in question should first of all, apply the common configuration to all nodes.
Second of all, it should install MicroK8s on all nodes.
After that, it needs to configure the master node appropriately.
Finally, it configures the other nodes to join the cluster as worker nodes so that they can execute our services.

There are more details related to Kubernetes such as node tainting and node selectors that won't be mentioned  in this report due to time constraints.

Also note that we won't get into too much details about how Ansible works in this report also due to time constraints and since it will become very lengthy. We'll only mention that the connection is established using SSH.

\begin{figure}[H]
    \centering
    \makebox[\textwidth]{\includegraphics[width=13cm]{src/assets/images/ansible-main-script.JPG}}
    \caption{Extract from the Ansible playbook}
    \label{fig:extract-from-ansible-playbook}
\end{figure}

\begin{figure}[H]
    \centering
    \makebox[\textwidth]{\includegraphics[width=12cm]{src/assets/images/ansible-hosts.JPG}}
    \caption{The Ansible hosts configuration}
    \label{fig:image-of-ansible-hosts}
\end{figure}
Note that we also use a hosts file for development where we test the setup extensively before trying it on production as seen from the figure just above.

% ----------------------- Application pipeline ----------------------- %
\section{Deployment process}
We will now explain the approach we adopted to continuously integrate, deploy and deliver our application as per the best practices of DevOps previously mentioned in \textbf{Chapter 3: State of the art}.
The deployment is currently supported on two different environments; the staging environment where we make sure our application works as intended, and of course the production environment.

\subsection{Adopted strategy}
Since we have multiple microservices, it would not be ideal to interact with our Kubernetes cluster from all of them whenever we make a change.
For that reason we use the following steps to deploy our application.
\begin{itemize}
    \item On code changes on any microservice, we create a new Docker image that we push to a custom container registry.
    \item Once the build is finished, we trigger a multi-project pipeline to a central repository we will call ilg-controller.
    \item The central repository will interact with our Kubernetes cluster to make the final deployment to the correct environment using a custom Helm package that we create.
\end{itemize}

This is further explained in detail in the following sub sections.

\subsection{Linking the cluster to GitLab}
Since we have a Kubernetes cluster and we also use GitLab, we will link the two using the GitLab Kubernetes Agent, or KAS for short.

\begin{figure}[H]
    \centering
    \makebox[\textwidth]{\includegraphics[width=15cm]{src/assets/images/gitlab-kubernetes-agent.JPG}}
    \caption{GitLab Kubernetes Agent}
    \label{fig:gitlab-kubernetes-agent}
\end{figure}

% CI Continuous Integration
\subsection{Continuous Integration}
When we as developers make code changes to the main branch of any microservice, tests are automatically executed.
If the tests are validated, a new Docker image that incorporates the changes we make will be pushed to the container registry and our deployment pipeline will then be triggered.
In other words, the new code will be integrated and that actually marks the end of this step.
\begin{figure}[H]
    \centering
    \makebox[\textwidth]{\includegraphics[width=10cm]{src/assets/images/ci-pipeline-extract.JPG}}
    \caption{Extract from the CI pipeline}
    \label{fig:extract-from-ci-pipeline}
\end{figure}

% CD Continuous Deployment
\subsection{Continuous Deployment}
Our ilg-controller repository (where the agent is configured) chooses the correct microservice to deploy depending on the trigger input.
It then sets the necessary environment variables for it and installs it using a custom Helm package we've created for the deployment.

\begin{figure}[H]
    \centering
    \makebox[\textwidth]{\includegraphics[width=11cm]{src/assets/images/cd-makefile-extract.JPG}}
    \caption{Extract from the deployment Makefile}
    \label{fig:extract-from-cd-makefile}
\end{figure}

% NOTE: Continuous delivery is a partly manual process where developers can deploy any changes to customers by simply clicking a button, while continuous deployment emphasizes automating the entirety of the process.

% Custom Helm Package
\subsection{Custom Helm Package}
Since we agreed to use Helm (refer to \textbf{Chapter 2: State of the art}), we've created a package -or in terms of Helm- a chart of our own to manage the Kubernetes cluster resources.

\begin{figure}[H]
    \centering
    \makebox[\textwidth]{\includegraphics[width=15cm]{src/assets/images/helm-chart-code-extract.JPG}}
    \caption{Extract from the custom Helm Chart}
    \label{fig:extract-from-custom-helm-chart}
\end{figure}

We even integrate a CI/CD pipeline for the Helm chart itself so that new packages are built automatically whenever we make changes to the manifest template files.

\begin{figure}[H]
    \centering
    \makebox[\textwidth]{\includegraphics[width=15cm]{src/assets/images/build-pacakge-stage.JPG}}
    \caption{Extract from the custom Helm package's pipeline}
    \label{fig:build-package-stage}
\end{figure}

\begin{figure}[H]
    \centering
    \makebox[\textwidth]{\includegraphics[width=15cm]{src/assets/images/gitlab-package-registry.JPG}}
    \caption{GitLab's package registry}
    \label{fig:gitlab-package-registry}
\end{figure}
\newpage

% TODO: Our deployed application
\subsection{Our deployed application}
With the previous work here done, we can now view our deployed application at https://lead-generator.incedo.net
and here is how everything looks like

\begin{figure}[H]
    \centering
    \makebox[\textwidth]{\includegraphics[width=15cm]{src/assets/app-screenshots/login.png}}
    \caption{Login page}
    \label{fig:login-page}
\end{figure}
\newpage
\begin{figure}[H]
    \centering
    \makebox[\textwidth]{\includegraphics[width=15cm]{src/assets/app-screenshots/campaigns.png}}
    \caption{Campaigns Overview page}
    \label{fig:campaigns-overview}
\end{figure}
\begin{figure}[H]
    \centering
    \makebox[\textwidth]{\includegraphics[width=15cm]{src/assets/app-screenshots/campaigns-add.png}}
    \caption{Add Campaign form}
    \label{fig:campaigns-add-form}
\end{figure}
\newpage
\begin{figure}[H]
    \centering
    \makebox[\textwidth]{\includegraphics[width=15cm]{src/assets/app-screenshots/campaigns-edit.png}}
    \caption{Edit Campaign form}
    \label{fig:campaigns-edit-form-1}
\end{figure}
\begin{figure}[H]
    \centering
    \makebox[\textwidth]{\includegraphics[width=15cm]{src/assets/app-screenshots/users.png}}
    \caption{View Users page}
    \label{fig:users-view-page-1}
\end{figure}
\newpage
\newpage
\begin{figure}[H]
    \centering
    \makebox[\textwidth]{\includegraphics[width=15cm]{src/assets/app-screenshots/campaigns-edit.png}}
    \caption{Edit Campaign form}
    \label{fig:campaigns-edit-form-2}
\end{figure}
\begin{figure}[H]
    \centering
    \makebox[\textwidth]{\includegraphics[width=15cm]{src/assets/app-screenshots/users.png}}
    \caption{View Users page}
    \label{fig:users-view-page-2}
\end{figure}\newpage
\begin{figure}[H]
    \centering
    \makebox[\textwidth]{\includegraphics[width=15cm]{src/assets/app-screenshots/users-create.png}}
    \caption{Add User form}
    \label{fig:users-add-form}
\end{figure}
\begin{figure}[H]
    \centering
    \makebox[\textwidth]{\includegraphics[width=15cm]{src/assets/app-screenshots/users-add-credit.png}}
    \caption{Users Add credit form}
    \label{fig:users-add-credit-form}
\end{figure}
\begin{figure}[H]
    \centering
    \makebox[\textwidth]{\includegraphics[width=15cm]{src/assets/app-screenshots/records.png}}
    \caption{Campaigns Records page}
    \label{fig:records-page}
\end{figure}
\begin{figure}[H]
    \centering
    \makebox[\textwidth]{\includegraphics[width=15cm]{src/assets/app-screenshots/analytics-1.png}}
    \caption{Campaigns Analytics page}
    \label{fig:campaigns-analytics-page-1}
\end{figure}
\begin{figure}[H]
    \centering
    \makebox[\textwidth]{\includegraphics[width=15cm]{src/assets/app-screenshots/analytics-2.png}}
    \caption{Campaigns Analytics page}
    \label{fig:campaign-analytics-page-2}
\end{figure}
\newpage

% ----------------------- Monitoring ----------------------- %
\section{Monitoring}
In order to make the most out of our Kubernetes setup, we also setup monitoring for our whole application stack.
The things we monitor vary from the cluster nodes' CPU and memory usage, to the logs of our microservices.

In order to get started, we setup the necessary manifest resources to expose the dashboard provided by Kubernetes.

\subsection{Monitoring the cluster nodes}
By using and configuring the Kubernetes community's Helm packages, we can very easily deploy what's called the \textbf{Kubernetes Prometheus Stack} which will give us access to the Grafana dashboard and to the Prometheus UI dashboard.

\subsubsection*{\underline{Prometheus}}
Prometheus is an open-source software used for event monitoring and alerting.
The way we do that is by collecting different kind of metrics as well as configuring the automatic sending of alerts in case of critical errors.
\begin{figure}[H]
    \centering
    \makebox[\textwidth]{\includegraphics[width=15cm]{src/assets/images/prometheus-metrics.JPG}}
    \caption{Viewing the metrics endpoints with Prometheus UI}
    \label{fig:prometheus-ui-metrics-endpoint}
\end{figure}

\subsubsection*{\underline{Grafana}}
Grafana is an open-source interactive visualization web application.
We use it here to monitor our cluster nodes, and most importantly, the metrics of our application collected by Prometheus.
\begin{figure}[H]
    \centering
    \makebox[\textwidth]{\includegraphics[width=15cm]{src/assets/images/grafana-nodes-monitoring.JPG}}
    \caption{Monitoring the cluster nodes with Grafana}
    \label{fig:grafana-node-monitoring}
\end{figure}

\subsection{Monitoring the logs}
In order to monitor the logs of our different microservices as well as any other applications running on the cluster efficiently, we need to

\begin{itemize}
    \item Collect the logs from the different Pods in the cluster
    \item Process and parse the logs to a uniform format, often json
    \item Send the logs to an indexing engine for easy querying
    \item We may also optionally create dashboards for our logs
\end{itemize}
The best way to go about doing this is by using the EFK Stack or more expressively the Elasticsearch, Fluentd and Kibana stack.

\subsubsection*{\underline{Elasticsearch \& Kibana}}
Elasticsearch is many things: it is often referred to as search engine, analytics engine or even database and all of that is true.
It can process a huge amount of logs with relatively minimal resources and is therefore the optimal choice for our cluster.

Kibana is a data visualization dashboard exclusively used with Elasticsearch.
It is always almost used whenever there is an Elasticsearch instance involved.

By using the open-source \textbf{Elastic Stack Kubernetes Helm Charts}, we can create our Elasticsearch and Kibana setup in the cluster.
\begin{figure}[H]
    \centering
    \makebox[\textwidth]{\includegraphics[width=15cm]{src/assets/images/kibana-home.JPG}}
    \caption{Kibana dashboard home screen}
    \label{fig:kibana-dashboard-home}
\end{figure}

\subsubsection*{\underline{Fluentd}}
Fluentd on the other hand is a data collector.
It runs in the background with the sole purpose of delivering the logs of all of our applications to us in a uniform way we specify.
We configure it to get the logs from the standard output of our Pods, parse them to json and send them to Elasticsearch.

We deploy Fluentd the manifest way, as opposed to using Helm Charts.
\begin{figure}[H]
    \centering
    \makebox[\textwidth]{\includegraphics[width=15cm]{src/assets/images/fluentd-config-snippet.JPG}}
    \caption{Fluentd configuration extract}
    \label{fig:fluentd-config-snippet}
\end{figure}

And now we can finally use some of the advanced features of Elasticsearch such as the filtering and sorting of our logs from the Kibana dashboard.

\begin{figure}[H]
    \centering
    \makebox[\textwidth]{\includegraphics[width=15cm]{src/assets/images/efk-stack-logs.JPG}}
    \caption{Indexing the logs in Kibana}
    \label{fig:kibana-logs-indexing}
\end{figure}

\newpage

% ----------------------- Technologies ----------------------- %
\section{Technologies}
This part is reserved for the presentation of all of the software used in the realization of the project and includes but is not limited to; programming languages, frameworks, technologies, etc...

\medskip
For a comparative analysis on some of our choices, see \textbf {Chapter 2: State of the art}.
% NOTE: Transform into a table if total page number goes over 70 or so

\begin{itemize}
    \item \textbf{Git:} \newline
          \begin{minipage}{\linewidth}
              \centering
              \includegraphics[width=2.5cm]{src/assets/logos/git_512x512.png}
              \captionof{figure}{Logo of Git}
          \end{minipage}
    \item \textbf{Docker:} \newline
          \begin{minipage}{\linewidth}
              \centering
              \includegraphics[width=4cm]{src/assets/logos/docker_512x512.png}
              \captionof{figure}{Logo of Docker}
          \end{minipage}
    \item \textbf{Playwright:} \newline
          \begin{minipage}{\linewidth}
              \centering
              \includegraphics[width=4.5cm]{src/assets/logos/playwright_512x512.png}
              \captionof{figure}{Logo of Playwright}
          \end{minipage}

          \newpage
    \item \textbf{NestJS:} \newline \cite{nestjs} A framework for building efficient, scalable Node.js web applications. \newline
          \begin{minipage}{\linewidth}
              \centering
              \includegraphics[width=3.7cm]{src/assets/logos/nestjs_512x512.png}
              \captionof{figure}{Logo of NestJS}
          \end{minipage}
    \item \textbf{ReactJS:} \newline A front-end JavaScript library that's often considered as a framework. \newline
          \begin{minipage}{\linewidth}
              \centering
              \includegraphics[width=3.7cm]{src/assets/logos/react_512x512.png}
              \captionof{figure}{Logo of ReactJS}
          \end{minipage}
    \item \textbf{PostgreSQL:} \newline \newline
          \begin{minipage}{\linewidth}
              \centering
              \includegraphics[width=8cm]{src/assets/logos/postgresql.png}
              \captionof{figure}{Logo of PostgreSQL}
          \end{minipage}

          \newpage
    \item \textbf{Sequelize:} \newline \cite{sequelize} A modern TypeScript and Node.js ORM for SQL databases. \newline
          \begin{minipage}{\linewidth}
              \centering
              \includegraphics[width=5cm]{src/assets/logos/sequelize_512x512.png}
              \captionof{figure}{Logo of Sequelize}
          \end{minipage}
    \item \textbf{GraphQL:} \newline \cite{graphql} A query language for APIs and a runtime for fulfilling those queries with existing data. \newline
          \begin{minipage}{\linewidth}
              \centering
              \includegraphics[width=3.7cm]{src/assets/logos/graphql_512x512.png}
              \captionof{figure}{Logo of GraphQL}
          \end{minipage}
    \item \textbf{GitLab:} \newline A DevOps software that allows secure collaboration and operations in a single application. \newline
          \begin{minipage}{\linewidth}
              \centering
              \includegraphics[width=3.5cm]{src/assets/logos/gitlab_512x512.png}
              \captionof{figure}{Logo of GitLab}
          \end{minipage}

          \newpage
    \item \textbf{Nginx:} \newline \cite{nginx} A multi-purpose web server can also be used as reverse proxy, load balancer, mail proxy and HTTP cache. \newline \newline
          \begin{minipage}{\linewidth}
              \centering
              \includegraphics[width=3.4cm]{src/assets/logos/nginx_512x512.png}
              \captionof{figure}{Logo of Nginx}
          \end{minipage}
    \item \textbf{Renovate:} \newline A software that provides automatic dependency updates with support for multiple languages. \newline
          \begin{minipage}{\linewidth}
              \centering
              \includegraphics[width=3cm]{src/assets/logos/renovate_200x200.png}
              \captionof{figure}{Logo of Renovate}
          \end{minipage}
    \item \textbf{Docker compose:} \newline A helper tool to run applications with multiple Docker containers using a yaml format. Often used in development. \newline \newline
          \begin{minipage}{\linewidth}
              \centering
              \includegraphics[width=4cm]{src/assets/logos/docker-compose_667x667.png}
              \captionof{figure}{Logo of Docker Compose}
          \end{minipage}

          \newpage
    \item \textbf{Visual Studio Code:} \newline A code editor that's used as an entire development environment with the help of extensions. \newline
          \begin{minipage}{\linewidth}
              \centering
              \includegraphics[width=2.5cm]{src/assets/logos/vscode_512x512.png}
              \captionof{figure}{Logo of VSCode}
          \end{minipage}
    \item \textbf{Linux:} \newline A Unix-like operating system for common daily use and more commonly for servers. \newline
          \begin{minipage}{\linewidth}
              \centering
              \includegraphics[width=3cm]{src/assets/logos/linux_512x512.png}
              \captionof{figure}{Logo of Linux}
          \end{minipage}
    \item \textbf{Make:} \newline GNU Make is used extensively throughout the application to save useful commands for the ease of maintainability. \newline
          \begin{minipage}{\linewidth}
              \centering
              \includegraphics[width=3cm]{src/assets/logos/makefile_512x512.png}
              \captionof{figure}{Logo of Makefile}
          \end{minipage}

          \newpage
    \item \textbf{MicroK8s:} \newline \cite{microk8s} MicroK8s is a simple production-grade conformant K8s. It is Lightweight and focused. \newline \newline
          \begin{minipage}{\linewidth}
              \centering
              \includegraphics[width=3.8cm]{src/assets/logos/microk8s.png}
              \captionof{figure}{Logo of MicroK8s}
          \end{minipage}
    \item \textbf{Kubernetes:} \newline The most powerful tool for managing containerized workloads in the cloud. \newline
          \begin{minipage}{\linewidth}
              \centering
              \includegraphics[width=3.5cm]{src/assets/logos/kubernetes_512x512.png}
              \captionof{figure}{Logo of Kubernetes}
          \end{minipage}
    \item \textbf{Ansible:} \newline An open-source software for provisioning, configuring and managing infrastructure. \newline
          \begin{minipage}{\linewidth}
              \centering
              \includegraphics[width=3.5cm]{src/assets/logos/ansible_200x200.png}
              \captionof{figure}{Logo of Ansible}
          \end{minipage}

          \newpage
    \item \textbf{Squid Proxy:} \newline \cite{squid} Squid is a caching proxy for the Web that supports HTTP, HTTPS, FTP and more. \newline
          \begin{minipage}{\linewidth}
              \centering
              \includegraphics[width=6cm]{src/assets/logos/squid-proxy.png}
              \captionof{figure}{Logo of Squid Proxy}
          \end{minipage}
    \item \textbf{GitLab agent for Kubernetes:} \newline \cite{kas-blog} A secure and reliable way to attach a Kubernetes cluster to GitLab. \newline
          \begin{minipage}{\linewidth}
              \centering
              \includegraphics[width=3.8cm]{src/assets/logos/gitlab-kubernetes-agent_512x512.png}
              \captionof{figure}{Logo of GitLab Agent for Kubernetes}
          \end{minipage}
    \item \textbf{Helm:} \newline Helm is the most popular package manager for Kubernetes. \newline
          \begin{minipage}{\linewidth}
              \centering
              \includegraphics[width=4cm]{src/assets/logos/helm_512x512.png}
              \captionof{figure}{Logo of Helm}
          \end{minipage}

          \newpage
    \item \textbf{Prometheus:} \newline Prometheus is a software application used for monitoring and alerting. \newline
          \begin{minipage}{\linewidth}
              \centering
              \includegraphics[width=4cm]{src/assets/logos/prometheus_512x512.png}
              \captionof{figure}{Logo of Prometheus}
          \end{minipage}
    \item \textbf{Grafana:} \newline Grafana is an open source solution for monitoring and analytics. \newline \newline
          \begin{minipage}{\linewidth}
              \centering
              \includegraphics[width=3.5cm]{src/assets/logos/grafana.png}
              \captionof{figure}{Logo of Grafana}
          \end{minipage}
    \item \textbf{LaTeX:} \newline \cite{latex-project} A high-quality document preparation and typesetting system for technical grade documents. \newline \newline
          \begin{minipage}{\linewidth}
              \centering
              \includegraphics[width=4cm]{src/assets/logos/latex_200x200.png}
              \captionof{figure}{Logo of The LaTeX Project}
          \end{minipage}

          \newpage
    \item \textbf{Elasticsearch:} \newline Elasticsearch is a search engine based on the Lucene library. Is it often used as part of a stack alongside Kibana and other tools. \newline
          \begin{minipage}{\linewidth}
              \centering
              \includegraphics[width=7cm]{src/assets/logos/elasticsearch-logo.png}
              \captionof{figure}{Logo of elasticsearch}
          \end{minipage}
    \item \textbf{Kibana:} \newline \cite{kibana} Kibana is an free and open frontend application that sits on top of the Elastic Stack, providing search and data visualization capabilities for data indexed in Elasticsearch. \newline
          \begin{minipage}{\linewidth}
              \centering
              \includegraphics[width=6.8cm]{src/assets/logos/kibana-logo.jpg}
              \captionof{figure}{Logo of Kibana}
          \end{minipage}
    \item \textbf{Fluentd:} \newline \cite{fluentd} Fluentd is an open source data collector, which lets you unify the data collection and consumption for a better use and understanding of data. \newline
          \begin{minipage}{\linewidth}
              \centering
              \includegraphics[width=4cm]{src/assets/logos/fluentd-logo.png}
              \captionof{figure}{Logo of Fluentd}
          \end{minipage}

          \newpage
    \item \textbf{Microsoft Teams:} \newline Microsoft Teams is a proprietary business communication platform developed by Microsoft. We use it to collaborate and share information. \newline
          \begin{minipage}{\linewidth}
              \centering
              \includegraphics[width=3.2cm]{src/assets/logos/ms-teams-logo.png}
              \captionof{figure}{Logo of Microsoft Teams}
          \end{minipage}
    \item \textbf{Asana:} \newline Asana is a work management platform available on web and and mobile. It is used internally by Incedo for various work and project management. \newline
          \begin{minipage}{\linewidth}
              \centering
              \includegraphics[width=4.5cm]{src/assets/logos/asana-logo.png}
              \captionof{figure}{Logo of Asana}
          \end{minipage}
    \item \textbf{Signal:} \newline Signal is a cross-platform centralized e2e encrypted instant messaging service. It is used by the team to share sensitive information. \newline
          \begin{minipage}{\linewidth}
              \centering
              \includegraphics[width=6cm]{src/assets/logos/signal-logo.png}
              \captionof{figure}{Logo of Signal}
          \end{minipage}

          \newpage
\end{itemize}

% ----------------------- Difficulties encountered ----------------------- %
\section{Difficulties encountered}
\subsection{Infrastructure as code}
This not per say a technical difficulty as it is more of a limitation from the side of Ionos cloudpanel; the service of IONOS Cloud that we use the provision the needed resources.
To have a full DevOps approach, even the infrastructure should be managed in code -also known as \glsxtrfull{iac}- using tools such as Terraform.
However, Ionos cloudpanel does not provide that functionality.
Therefore we had to scrape off the idea and simply create the virtual machines from the UI every time we need them.

\subsection{Stateful Kubernetes}
Another issue we had is data persistence in Kubernetes.
Kubernetes is -by design- used for stateless applications.
However, since we're using an Elasticsearch database, we definitely needed a way to persist our data across system reboots or crashes.

First of all, we used IONOS -our cloud provider- to create the physical Block Storage in question that we used as the data source.
We started with an SSD of 40 Gigabytes then scaled it up to 80 Gigs.

\begin{figure}[H]
    \centering
    \makebox[\textwidth]{\includegraphics[width=15cm]{src/assets/images/ionos-block-storage.JPG}}
    \caption{Creating a block storage in IONOS}
    \label{fig:ionos-block-storage}
\end{figure}

Secondly, we had to logically mount the previously created block storage to a directory such as \textbf{/mnt/block} inside the controller virtual machine.

From there, we create a self-managed NFS Server on the host system that we export to the private network or in other words to all of the cluster nodes.

Now, for the stateful cluster components to use the NFS Server, we needed to setup an NFS client provisioner inside the Kubernetes cluster itself. This was as simple as using the \textbf{NFS Subdir External Provisioner} Helm Chart.

Lastly, our cluster components are now allowed to provision persistent and dynamic storage resources whenever they need it.

\setcounter{secnumdepth}{0} % Set the section counter to 0 so next section is not counted in toc
% ----------------------- Conclusion ----------------------- %
\section{Conclusion}
A successful project comes from a successful development environment as well as a clean production environment with a good and well thought of monitoring solutions.
For this reason, we spent quite a good amount of time setting up the infrastructure in a clean and scalable way using various tools and scripts which all conform to the modern ways of how DevOps should be done.
